\section{Introduction}
\label{sec:introduction}

Is generative AI reshaping employment patterns? A growing literature investigates a post-2022 age gradient in occupations most exposed to artificial intelligence: employment of young workers declines relative to older workers, while no such gradient appears in less-exposed occupations. \cite{brynjolfsson2025} find this pattern in US payroll data, \cite{lodefalk2026} document a similar gradient in Sweden, and \cite{teeselink2025} reports consistent UK evidence. On the other hand, \cite{humlum2026} and \cite{kauhanen2026} report null patterns from Denmark and Finland.

We document Norwegian employment patterns. Using population-level register data from the Norwegian employer reporting system for 2021 to 2026, we construct monthly employment series at the level of cells defined by occupation, age group, sector, and AI exposure quintile, indexed to October 2022, the month preceding the November 30, 2022 release of ChatGPT. We rank occupations by the Eloundou et al.\ GPT-4 exposure measure. Selected high-exposure occupations show falling employment among the youngest workers, but across the full exposure distribution the most-exposed quintile does not decline relative to the least-exposed for workers under 40. The negative exposure gradient appears instead among workers aged 41 to 50, and is concentrated in the private sector. Public-sector employment shows no corresponding gradient. Cash earnings within occupation-age cells show no systematic divergence by exposure quintile.

The data cover the entire employed population at monthly frequency, with age assigned at monthly precision. Our main exposure measure is the theoretical task measure of \citet{eloundou2024}; the appendix reports the automation and augmentation shares of \citet{handa2025}, based on revealed Claude usage, as an alternative.

Norway is also a setting where AI tools have diffused rapidly. A 2026 Microsoft report ranks Norway third globally in population-level AI usage, after the UAE and Singapore \citep{microsoft2026ai}. OECD data for 2025 place Norway first among 33 European countries for AI use at work \citep{oecd2026ai}. At the enterprise level, the share of Norwegian firms (10 or more employees) reporting AI use rose from 11 percent in 2021 to 30 percent in 2025 \citep{ssb2025ikt}, placing Norway well above the EU average of 20 percent; the Nordic countries as a group lead European enterprise adoption \citep{bick2026}. Internet access (99 percent) and OECD top-five broadband density provide the infrastructure for rapid diffusion \citep{oecd2024digitalnorway}. If AI-related labor market effects are present, a high-adoption economy is where they would appear first.
